\section{Verdikt}
\label{sec:verdikt-rdy}

Rot gesetzt: eine Empfehlung unter den Annahmen aus
Abschnitt~\ref{sec:annahmen-rdy}, kein Befund.

\subsection{Votum}

\vd{Meiden zum Kurs von \USDje{\RdyKurs}
(\INRje{\RdyKursRupien} je Aktie), und zwar nicht, weil der Preis zu hoch
w\"are, sondern weil er nicht pr\"ufbar ist. Eine Einstiegsschwelle l\"asst
sich f\"ur diesen Titel nicht angeben, ohne die eine Gr\"o{\ss}e zu
sch\"atzen, die das Unternehmen nicht ver\"offentlicht. Der Horizont dieses
Votums ist entsprechend kurz: Es gilt, bis ein Bericht den Beitrag des
auslaufenden Gesch\"afts beziffert oder zwei weitere Quartale die neue
Ertragsbasis zeigen.}

\vd{F\"ur den Halter: nicht aufstocken; die Position ist eine Wette auf eine
unbekannte Gr\"o{\ss}e, und sie wird durch Zukaufen nicht besser. Ein Verkauf
dr\"angt sich nicht auf, solange die Nettokasse steht. F\"ur den
Nichthalter: nicht einsteigen. Es gibt keinen Preis, den dieses Dokument
nennen kann, ohne zu raten.}

\subsection{Begr\"undung aus den Befunden}

\vd{\emph{Was dagegen spricht.} Der Ertrag f\"allt seit zwei Perioden, und
zwar beschleunigt: \PctBetrag{\RdyErgebnisWachstum} im Gesch\"aftsjahr,
\PctBetrag{\RdyQuartalGewinnWachstum} im Quartal danach. Die Rohertragsmarge
sank um mehr als f\"unf Prozentpunkte. Der operative Mittelzufluss des
j\"ungsten Quartals war negativ. \"Uber zehn Jahre erwirtschaftete der Titel
\Pct{\RdyRenditeZehn} pro Jahr gegen \Pct{\Huerde} der Branche. Das ist
keine Momentaufnahme, sondern der Befund eines Jahrzehnts.}

\vd{\emph{Was daf\"ur spricht.} Eine Nettokasse von
\MioINR{\RdyNettokasse}, eine Eigenkapitalquote von
\Pct{\RdyEigenkapitalquote} und ein Kurs-Buchwert-Verh\"altnis von
\Faktor{\RdyKursBuchwert}, weniger als die H\"alfte dessen, was
AstraZeneca kostet. Das geforderte Wachstum bei konstantem Vielfachen
betr\"agt \Pct{\RdyWachstumVielfachZehn} pro Jahr, und \"uber neun Jahre hat
das Unternehmen \Pct{\RdyCagr} geliefert. Wer diese Reihe f\"ur
aussagekr\"aftig h\"alt, findet den Titel billig.}

\vd{\emph{Warum das Votum trotzdem meiden lautet.} Weil dieselbe Reihe, sieben
statt neun Jahre lang gerechnet, \Pct{\RdyCagrSieben} pro Jahr ergibt. Zwei
Startjahre, zwei Vorzeichen, dieselben belegten Zahlen. Eine Bewertung, die
davon abh\"angt, welches Jahr man f\"ur den Anfang h\"alt, ist keine
Bewertung. Bei einer Rate \"uber der H\"urde gibt es \"uberhaupt keine
Preisobergrenze mehr; die Rechnung h\"ort dann auf, eine Frage zu
beantworten. Bei der niedrigeren Rate liegt die Schwelle bei
\INRje{\RdySchwelleHistorieSieben} gegen einen Kurs von
\INRje{\RdyKursRupien}. Zwischen \enquote{jeder Preis ist gerechtfertigt} und
\enquote{ein Sechstel des heutigen} liegt kein Urteil, sondern eine
Wissensl\"ucke.}

\vd{\emph{Wo dieses Votum vom Konsens abweicht.} Im Niveau, nicht in der
Richtung. \Stueck{\RdyKonsensHaeuser} H\"auser sagen \enquote{Hold}, dieses
Dokument sagt meiden; das ist derselbe Befund in sch\"arferer Form. Das
mittlere Kursziel von \USDje{\RdyKonsensZiel}
(\INRje{\RdyKonsensZielRupien}) liegt \Pct{\RdyKonsensUeberVariante} \"uber
der mittleren Schwelle dieses Berichts. Die Abdeckung ist mit
\Stueck{\RdyKonsensHaeuser} Sch\"atzungen d\"unn, und eine d\"unne Abdeckung
ist ein schw\"acheres Gegenargument als eine dichte.}

\subsection{Einstiegsschwelle}
\label{sec:schwelle-rdy}

\begin{table}[htbp]
\centering
\caption{Einstiegsschwelle auf zehn Jahre bei einer H\"urde von
\Pct{\Huerde}, je Endwertannahme. Kurs heute: \INRje{\RdyKursRupien} je
Aktie.}
\label{tab:rdy-schwellen-drei}
\begin{tabular}{llr}
\toprule
Endwertannahme & unterstelltes Wachstum & Schwelle\\
\midrule
Buchwert je Aktie       & keines & \INRje{\RdySchwelleZehn}\\
Kurs von heute, nominal & keines & \INRje{\RdySchwelleVarZehn}\\
konstantes Vielfaches   & eigene Rate FY2019--FY2026 & \INRje{\RdySchwelleHistorieSieben}\\
konstantes Vielfaches   & eigene Rate FY2018--FY2026 & keine Obergrenze\\
\bottomrule
\end{tabular}
\end{table}

\vd{Die vierte Zeile ist der Grund des Votums. Bei einer unterstellten Rate
oberhalb der H\"urde \"ubersteigt der Endwert allein den Einstiegspreis, und
jeder Preis erreicht die H\"urde; die Rechnung liefert dann keine Grenze mehr.
Das ist kein Rechenfehler, sondern die Aussage, dass die Frage bei dieser
Annahme keinen Inhalt mehr hat. Eine Zahl, die zwischen \enquote{unbegrenzt}
und \INRje{\RdySchwelleHistorieSieben} liegt, w\"are eine Erfindung.}

\FloatBarrier

\subsection{Ausl\"oser, die das Votum drehen}

\vd{\emph{Zum Kauf.} Erstens: ein Bericht, der den Umsatz- oder
Ergebnisbeitrag des Lenalidomid-Gesch\"afts beziffert. Das w\"are die
Ver\"offentlichung der fehlenden Gr\"o{\ss}e, und sie macht den Titel
bewertbar, und zwar in beide Richtungen. Zweitens: zwei aufeinanderfolgende Quartale
mit stabilem oder steigendem Rohertrag im Segment Global Generics; das w\"are
der Beleg, dass der Boden gefunden ist. Drittens: die US-Zulassung eines der
drei anh\"angigen Biosimilars, insbesondere Rituximab.}

\vd{\emph{Zum endg\"ultigen Ausscheiden.} Ein weiteres Quartal mit negativem
operativem Mittelzufluss, oder eine Beanstandung der US-Beh\"orde h\"oherer
Stufe als \enquote{Voluntary Action Indicated} an einem der
Formulierungswerke.}

\subsection{Was dieses Votum nicht wissen kann}

\vd{Die Gr\"o{\ss}e, an der alles h\"angt, ist genannt und nicht beziffert.
Alles Weitere folgt daraus.}

\vd{Drei Argumente gegen das eigene Votum.}

\vd{Erstens: \enquote{Meiden} ist ein bequemes Votum, wenn eine Angabe fehlt.
Es kostet nichts und kann nie falsch aussehen. Ein Anleger, der die Reihe
FY2018--FY2026 f\"ur die aussagekr\"aftigere h\"alt (und sie ist die
l\"angere), kauft hier ein Unternehmen mit Nettokasse zum
\Faktor{\RdyKursBuchwert}-fachen Buchwert nach einem R\"uckgang von
\PctBetrag{\RdyAbstandHoch}. Das ist eine verteidigbare Gegenposition.}

\vd{Zweitens: Der Zukauf des Nikotinersatzgesch\"afts wird hier als
margenschw\"acherer Ersatz behandelt. Belegt ist das nicht. Belegt ist nur,
dass die Segmentmarge gefallen ist, w\"ahrend dieses Gesch\"aft hinzukam;
eine eigene Marge des zugekauften Portfolios weist der Bericht nicht aus.
Die Zuschreibung ist eine Einsch\"atzung und steht deshalb blau.}

\vd{Drittens: Das Quartal zum 30.~Juni ist bei diesem Unternehmen
traditionell das schw\"achste des Gesch\"aftsjahres, und ein Quartal ist
keine Reihe. Wer dem Jahresabschluss mehr Gewicht gibt als dem Quartal, sieht
ein Unternehmen mit \INRje{\RdyFcfVoll} freiem Mittelzufluss je Aktie nahe
seinem eigenen Median, und nicht den Beginn eines Absturzes.}
